\documentclass[11pt,a4paper]{article}
\usepackage{fontspec,amsmath,amssymb,amsthm,microtype}
\setmainfont{lmroman10-regular.otf}[BoldFont=lmroman10-bold.otf,ItalicFont=lmroman10-italic.otf,BoldItalicFont=lmroman10-bolditalic.otf]
\usepackage{xcolor}
\usepackage[most]{tcolorbox}
\usepackage{mathrsfs}
\providecommand{\lt}{<}
\providecommand{\gt}{>}
\setmonofont{DejaVuSansMono.ttf}[BoldFont=DejaVuSansMono-Bold.ttf,ItalicFont=DejaVuSansMono-Oblique.ttf,BoldItalicFont=DejaVuSansMono-BoldOblique.ttf,Scale=0.86]
\newfontfamily\leanfallback{STIXTwoMath-Regular.otf}[Scale=0.86]
\definecolor{leanrule}{gray}{0.62}
\definecolor{leanback}{gray}{0.962}
\definecolor{leantitle}{gray}{0.22}
\definecolor{leanmuted}{gray}{0.38}
\newtcolorbox{leanbox}[1]{enhanced,breakable,colback=leanback,colframe=leanrule,
  boxrule=0.4pt,arc=0.6mm,left=2.4mm,right=2.4mm,top=1.8mm,bottom=1.8mm,
  before skip=0.8em,after skip=1em,fonttitle=\small\bfseries,coltitle=leantitle,
  colbacktitle=leanback,titlerule=0.3pt,title={#1}}
\newlength{\leancharwidth}
\newenvironment{leancode}{\par\smallskip\begingroup\ttfamily\footnotesize\raggedright
  \settowidth{\leancharwidth}{M}\setlength{\parskip}{0pt}\setlength{\parindent}{0pt}
  \frenchspacing}{\par\endgroup\smallskip}
\newcommand{\leanline}[2]{\hangindent=\dimexpr#1\leancharwidth+4\leancharwidth\relax
  \hangafter=1\noindent\hspace*{#1\leancharwidth}#2\par}
\newcommand{\leansource}[1]{\par{\footnotesize\color{leanmuted}#1}\par}
\newcommand{\leanlabel}[1]{\par\smallskip\noindent{\small\bfseries #1}\par}
\usepackage[margin=25mm]{geometry}
\usepackage[unicode,hidelinks]{hyperref}
\hypersetup{pdfauthor={LangevinLean Project}}
\setlength{\emergencystretch}{3em}
\setlength{\parskip}{0.4em}
\setlength{\parindent}{0pt}
\allowdisplaybreaks[2]
\raggedbottom
\title{\texorpdfstring{The integrated common-noise local error for OBABO}{The integrated common-noise local error for OBABO}}
\author{LangevinLean Project}
\date{}
\begin{document}\maketitle
Draft mathematical exposition; these notes may contain mistakes.


This section turns the \(L^{2}\) Hamiltonian local comparison into the integrated cost of the original finite-dimensional maps, retaining the same initial state and both noise variables on one probability space.



The local kernel discrepancy is a common-noise transport cost. It is obtained by integrating the deterministic Verlet error against a stationary Hamiltonian comparator in the contracting phase metric.


\section{\texorpdfstring{The explicit local-error bound}{The explicit local-error bound}}


The local-error bound is \(b_h=2h^2(\sqrt{Md}+\sqrt M\sqrt d+h\sqrt M\sqrt{Md}/4)\). Its nonnegativity follows directly from \(h\ge0\) and square-root nonnegativity. It is \(O(h^{2}\sqrt d)\) with coefficients depending only on \(M\) and the allowed step-size range. The following calculation proves that this expression bounds the numerical error.



The local-error bound retains the explicit step size, friction, smoothness, and dimension factors obtained from the Hamiltonian comparator. Its order in the step is quadratic. This order is needed because geometric accumulation will later divide by a contraction gap of first order.

\begin{leanbox}{Lean: \texttt{obaboLocalBudget}, \texttt{obaboLocalBudget\_\allowbreak{}nonneg}}
\leansource{\texttt{obaboLocalBudget} \quad KineticLangevin/\allowbreak{}OBABOLocalCost.lean, lines 19–20}
\begin{leancode}
\leanline{0}{def\ obaboLocalBudget\ (M\ h\ d\ :\ ℝ)\ :\ ℝ\ :=}
\leanline{2}{2\ *\ h\ \textasciicircum{}\ 2\ *\ (Real.sqrt\ (M\ *\ d)\ +\ Real.sqrt\ M\ *\ Real.sqrt\ d\ +\ Real.sqrt\ M\ *\ Real.sqrt\ (M\ *\ d)\ *\ h\ /\ 4)}
\end{leancode}
\leansource{\texttt{obaboLocalBudget\_\allowbreak{}nonneg} \quad KineticLangevin/\allowbreak{}OBABOLocalCost.lean, lines 22–24}
\begin{leancode}
\leanline{0}{theorem\ obaboLocalBudget\_nonneg\ (M\ h\ d\ :\ ℝ)\ (hh\ :\ 0\ ≤\ h)\ :\ 0\ ≤\ obaboLocalBudget\ M\ h\ d}
\end{leancode}
\leanlabel{How the Lean proof follows the argument}
\texttt{obabo\allowbreak{}Local\allowbreak{}Budget} names the mathematical object defined above. Its defining expression appears in the adjacent statement; subsequent identities may unfold this definition directly.




\texttt{obabo\allowbreak{}Local\allowbreak{}Budget\_\allowbreak{}nonneg} gives the formal counterpart of this step.
\end{leanbox}

\section{\texorpdfstring{The full common-noise defect before taking its cost}{The full common-noise defect before taking its cost}}


Construct the \(L^{2}\) Hamiltonian path at the first thermostat output. The strengthened pointwise Hamiltonian lifting theorem identifies its endpoint almost everywhere with the finite-dimensional exact Hamiltonian flow at that same initial state. The lifted OBABO theorem and final-thermostat representative theorem therefore identify the two outputs jointly, including both original noise variables. Subtracting the phase representatives and applying the weighted-transform \(L^{2}\) identity gives the exact integrated weighted cost of the original vector map discrepancy. The stationary-reference local estimate then proves that this cost is at most \(b_h\). The only reference-law input is the law of the first thermostat output; normalized target initialization and exact thermostat invariance verify that hypothesis in the full sampler theorem.



Both updates use exactly the same thermostat noises. The weighted random-variable discrepancy is therefore a coupling cost for their output measures. Integrating the previously derived deterministic local defect under the stationary reference yields the stated kernel error.

\begin{leanbox}{Lean: \texttt{lp\_\allowbreak{}obabo\_\allowbreak{}hamiltonian\_\allowbreak{}local\_\allowbreak{}cost}}
\leansource{\texttt{lp\_\allowbreak{}obabo\_\allowbreak{}hamiltonian\_\allowbreak{}local\_\allowbreak{}cost} \quad KineticLangevin/\allowbreak{}OBABOLocalCost.lean, lines 26–59}
\begin{leancode}
\leanline{0}{theorem\ lp\_obabo\_hamiltonian\_local\_cost}
\leanline{4}{(f\ :\ EuclideanSpace\ ℝ\ ι\ →\ ℝ)\ (m\ γ\ h\ δ\ σ\ :\ ℝ)\ (M\ :\ ℝ≥0)}
\leanline{4}{(hf\ :\ Differentiable\ ℝ\ f)\ (hsc\ :\ StrongConvexOn\ univ\ m\ f)\ (hLip\ :\ LipschitzWith\ M\ (gradient\ f))}
\leanline{4}{(hm\ :\ 0\ <\ m)\ (hγ\ :\ 0\ <\ γ)\ (hM\ :\ 0\ <\ (M\ :\ ℝ))\ (hh\ :\ 0\ ≤\ h)}
\leanline{4}{(hfriction\ :\ 4\ *\ (M\ :\ ℝ)\ ≤\ γ\ \textasciicircum{}\ 2)\ (hδ\ :\ |δ|\ ≤\ 1)}
\leanline{4}{(Z\ :\ Lp\ (EuclideanSpace\ ℝ\ ι)\ 2\ P\ ×\ Lp\ (EuclideanSpace\ ℝ\ ι)\ 2\ P)}
\leanline{4}{(Ξ\ Ξ'\ :\ Lp\ (EuclideanSpace\ ℝ\ ι)\ 2\ P)}
\leanline{4}{(hlaw\ :\ P.map\ (lpPhaseCoe\ P\ (ornsteinStep\ δ\ σ\ Ξ\ Z))\ =\ kineticGibbsTarget\ f)\ :}
\leanline{4}{eLpNorm\ (fun\ ω\ ↦\ weightedPhaseTransform\ (1\ /\ (M\ :\ ℝ))\ (1\ /\ γ)}
\leanline{6}{(obaboStep\ (gradient\ f)\ h\ δ\ σ\ (lpPhaseCoe\ P\ Z\ ω)\ (Ξ\ ω)\ (Ξ'\ ω)\ -}
\leanline{8}{ornsteinStep\ δ\ σ\ (Ξ'\ ω)}
\leanline{10}{(hamiltonianFlow\ f\ M\ hLip\ h\ hh\ (ornsteinStep\ δ\ σ\ (Ξ\ ω)\ (lpPhaseCoe\ P\ Z\ ω))\ h)))\ 2\ P\ ≤}
\leanline{6}{ENNReal.ofReal\ (obaboLocalBudget\ M\ h\ (Fintype.card\ ι\ :\ ℝ))}
\end{leancode}
\leanlabel{How the Lean proof follows the argument}
\texttt{lp\_\allowbreak{}obabo\_\allowbreak{}hamiltonian\_\allowbreak{}local\_\allowbreak{}cost} gives the formal counterpart of this step. The main intermediate facts are \texttt{A}, \texttt{N}, \texttt{R}, \texttt{hlocal}, \texttt{h\allowbreak{}N}, \texttt{h\allowbreak{}R} in that order. Named proof applications include \texttt{ENNReal.\allowbreak{}of\allowbreak{}Real\_\allowbreak{}le\_\allowbreak{}of\allowbreak{}Real}; the adjacent proof shows their hypotheses and order of use.
\end{leanbox}


Earlier results used here: \cite[Section~3.13.2]{LangevinLean2026ULASC}.

\begin{thebibliography}{ULA-LSI}
\bibitem[ULA-SC]{LangevinLean2026ULASC}
LangevinLean Project. \emph{Unadjusted Langevin under Strong Convexity}. LL-01, 2026. Draft report, content revision d34ead06987c. \url{https://odakyildiz.github.io/LangevinLean/docs/pdf/UnadjustedLangevin/FullProof.pdf}.
\end{thebibliography}

\end{document}
